\documentclass[10pt,letterpaper]{article}
\usepackage[margin=.67in,top=.56in,bottom=.54in]{geometry}
% Shared style for Saange Tamang's academic and outreach CVs.
% TeX Gyre Pagella ships with standard TeX distributions; no proprietary fonts.
\usepackage[T1]{fontenc}
\usepackage{tgpagella}
\usepackage{xcolor,enumitem,tabularx,microtype,fancyhdr}
\definecolor{ink}{HTML}{202124}
\definecolor{accent}{HTML}{4F685A}
\definecolor{muted}{HTML}{56545A}
\definecolor{rulegray}{HTML}{C9C5C6}
\usepackage[colorlinks=true,urlcolor=accent,linkcolor=accent,
  pdfauthor={Saange Tamang},pdfsubject={Wireless Communications and Machine Learning},
  pdfkeywords={wireless communications, machine learning, reinforcement learning, graph learning}]{hyperref}
\color{ink}
\setlength{\parindent}{0pt}
\setlength{\parskip}{2pt}
\setlength{\emergencystretch}{2em}
\setlist[itemize]{leftmargin=13pt,label=\textbullet,itemsep=3pt,topsep=3pt,parsep=0pt,partopsep=0pt}
\newcommand{\sect}[1]{\par\vspace{8pt}{\color{accent}\fontsize{11}{13}\selectfont\bfseries #1}\par\vspace{2pt}{\color{rulegray}\hrule height .35pt}\vspace{4pt}}
\newcommand{\entry}[2]{\textbf{#1}\hfill{\color{muted}\fontsize{10}{12}\selectfont #2}\par}
\newcommand{\meta}[1]{{\color{muted}\fontsize{10}{12}\selectfont #1}\par}
\newcommand{\paperlink}{\href{https://github.com/Sange-creator/graph-attention-double-dqn-cellular-handover/tree/main/paper}{Manuscript}}
\newcommand{\codelink}{\href{https://github.com/Sange-creator/graph-attention-double-dqn-cellular-handover}{Research code}}
\newcommand{\nameheader}{%
  {\fontsize{28}{31}\selectfont\bfseries Saange Tamang}\par
  \vspace{3pt}
  {\fontsize{10}{12}\selectfont Electronics, Communication \& Information Engineer\enspace |\enspace Licensed Engineer, Nepal Engineering Council}\par
  \vspace{5pt}
  {\fontsize{10}{12}\selectfont Nepal\enspace |\enspace\href{mailto:connectwithsaange@gmail.com}{connectwithsaange@gmail.com}\enspace |\enspace
  \href{https://github.com/Sange-creator}{GitHub}\enspace |\enspace\href{https://www.linkedin.com/in/saange-tamang/}{LinkedIn}}\par
}

\hypersetup{pdftitle={Saange Tamang - Research CV}}
\pagestyle{empty}
\renewcommand{\sect}[1]{%
  \par\vspace{6pt}{\color{accent}\fontsize{11}{13}\selectfont\bfseries #1}%
  \par\vspace{2pt}{\color{rulegray}\hrule height .35pt}\vspace{3pt}%
}
\setlist[itemize]{itemsep=2pt,topsep=3pt}
\begin{document}
\fontsize{10.5}{12.7}\selectfont
\nameheader

\sect{Education}
\entry{Tribhuvan University, Institute of Engineering}{2022--2026}
B.E. in Electronics, Communication and Information Engineering\par
\meta{Paschimanchal Campus, Pokhara, Nepal \enspace|\enspace 79.02\% aggregate, First Division}

\sect{Research experience}
\entry{Cellular handover with graph reinforcement learning}{2026}
\meta{First author: Saange Tamang. Co-author: Sulav Kandel.}
\begin{itemize}
\item Led the technical work: implemented the Python simulator and PyTorch controller, trained graph-attention and dueling Double DQN policies, and validated the evaluation pipeline.
\item Designed and ran a matched comparison of three frozen policies against a configured LTE A3-style baseline across 639 reserved Kathmandu-area cases and 540 Pokhara transfer groups.
\item In Kathmandu-area simulations, aggregate throughput rose by 8.86--9.99\% and handovers per hour fell by 36.1--53.0\%. The share of short completed stays increased for all three policies.
\item Co-wrote the paper with Sulav Kandel, who gathered field observations and checked site coordinates against OpenStreetMap.
\end{itemize}

\sect{Conference paper}
\textbf{Saange Tamang}, Sulav Kandel. \emph{Graph-Attention Double-DQN for Cellular Handover: A Paired Simulation Evaluation.}\par
\meta{Accepted for presentation at the 19th IOE Graduate Conference (2026).}
\paperlink\enspace |\enspace\codelink

\sect{Telecommunications experience}
\entry{Nepal Telecom (NTC), Pokhara --- Engineering Intern}{12 Jan--27 Mar 2026}
\begin{itemize}
\item Joined a vehicle-based cellular drive test and observed how engineers collected measurements along the route. Accompanied the team on a cell tower maintenance visit and examined installed routers and server-room infrastructure.
\item Explored hybrid mesh topology in Cisco Packet Tracer. Attended training led by senior engineers on cell-site connectivity, DNS, NAT, VPNs, and international upstream links.
\end{itemize}

\sect{Undergraduate projects}
\entry{LTE handover optimization and load balancing}{2025--2026}
\meta{Four-person team, semesters 7--8. Supervisor: Asst. Prof. Hom Nath Tiwari.}
Contributed to implementation and evaluation. The subsequent two-author conference study grew from this project and focused on handover optimization.\enspace
\href{https://github.com/Sange-creator/GNN-DQN-Handover-LoadBalancing}{Project code/report}

\vspace{3pt}
\entry{SolVote: Electronic Voting Prototype}{2025}
\meta{Undergraduate minor project, semester 6}
Helped implement and evaluate a team prototype for voter registration, Solana token issuance, and voting.\enspace
\href{https://github.com/Sange-creator/Project-Solvote}{Project code}

\sect{Technical preparation}
\textbf{Research:} Python, PyTorch, NumPy, Matplotlib, Git, LaTeX. Paired evaluation and cluster-bootstrap intervals.\par
\textbf{Networking:} LTE handover, RSRP/SINR, Cisco Packet Tracer, Wireshark.
\href{https://github.com/Sange-creator/Computer-Network}{Networking labs}\par
\textbf{Interests:} Cellular mobility and 5G/6G RAN. Graph learning, reinforcement learning, and model evaluation.

\end{document}
